\documentclass[11pt, a4paper]{article}

\usepackage[T1]{fontenc}
\usepackage[utf8]{inputenc}
\usepackage{tgpagella}
\usepackage[spanish, es-noshorthands]{babel}
\usepackage{amsmath, amssymb, bm}
\usepackage{tabularx}
\usepackage{booktabs}
\usepackage[most]{tcolorbox}
\usepackage[margin=2.5cm]{geometry}
\usepackage{ragged2e}
\usepackage{fancyhdr}
\usepackage{tikz}

\pagestyle{fancy}
\fancyhf{}
\setlength{\headheight}{16pt}
\lhead{\textbf{UCB Sede Tarija} | Ing. Mecatrónica}
\chead{}
\rhead{IMT-121 Circuitos Electrónicos I | Puente a EC3}
\lfoot{Prof: M.Sc. Ing. Bernardo Quiroga Turdera}
\rfoot{Página \thepage}
\renewcommand{\headrulewidth}{0.4pt}

\definecolor{ucbblue}{RGB}{0, 51, 102}

\begin{document}

\begin{tcolorbox}[colback=ucbblue!5!white, colframe=ucbblue, title=\textbf{Transición de Estado Estacionario a Régimen Transitorio}]
    \centering \textbf{Puente Conceptual: De EC2 a EC3}
\end{tcolorbox}

\section*{1. Los Límites de las Redes Resistivas}
Hasta ahora hemos estudiado circuitos bajo el régimen \textbf{DC Estático}. Si alteramos un voltaje, todo el circuito responde de forma instantánea gobernado puramente por álgebra ($V = IR$).

¿Qué ocurre físicamente cuando conmutamos un interruptor? ¿La tensión o la corriente cambian en cero segundos?
\begin{itemize}
    \item Los circuitos puramente resistivos \textbf{no almacenan energía}.
    \item \textbf{Los Capacitores ($C$) y los Inductores ($L$)} sí la almacenan. Su presencia introduce \textbf{el Tiempo ($t$)} en el circuito.
\end{itemize}

\section*{2. Relaciones Constitutivas y Continuidad (La Inercia Eléctrica)}

\vspace{0.2cm}
\noindent
\begin{minipage}[t]{0.48\textwidth}
    \textbf{\textcolor{ucbblue}{El Capacitor ($C$):}}
    \begin{itemize}
        \item Relación: $i_C(t) = C \frac{dv_C}{dt}$
        \item \textbf{Inercia de Tensión:} El voltaje a través de un capacitor ideal no puede cambiar instantáneamente.
        \item Regla: $\mathbf{v_C(0^+) = v_C(0^-)}$
    \end{itemize}
\end{minipage}
\hfill
\begin{minipage}[t]{0.48\textwidth}
    \textbf{\textcolor{ucbblue}{El Inductor ($L$):}}
    \begin{itemize}
        \item Relación: $v_L(t) = L \frac{di_L}{dt}$
        \item \textbf{Inercia de Corriente:} La corriente a través de un inductor ideal no puede cambiar instantáneamente.
        \item Regla: $\mathbf{i_L(0^+) = i_L(0^-)}$
    \end{itemize}
\end{minipage}

\vspace{0.4cm}
\section*{3. Evolución Temporal Cualitativa (Circuitos RC y RL)}
\begin{center}
\begin{tikzpicture}[scale=0.9]
    % RC
    \draw[thick, ->] (0,0) -- (4,0) node[anchor=north] {$t$};
    \draw[thick, ->] (0,0) -- (0,2.5) node[anchor=east] {$v_C(t)$};
    \draw[blue, thick] (0,0) .. controls (0.5,2) and (1.5,2) .. (3.5,2);
    \draw[dashed] (0,2) node[left] {$V_f$} -- (3.5,2);
    \node[anchor=north] at (1.5,-0.5) {\textbf{Carga Capacitor}};
    
    % RL
    \draw[thick, ->] (6,0) -- (10,0) node[anchor=north] {$t$};
    \draw[thick, ->] (6,0) -- (6,2.5) node[anchor=east] {$i_L(t)$};
    \draw[red, thick] (6,0) .. controls (6.5,2) and (7.5,2) .. (9.5,2);
    \draw[dashed] (6,2) node[left] {$I_f$} -- (9.5,2);
    \node[anchor=north] at (7.5,-0.5) {\textbf{Carga Inductor}};
\end{tikzpicture}
\end{center}

Observe que desde $t = 0^+$ hasta alcanzar el estado final, la variable experimenta un \textbf{transitorio} gradual dictado por una escala de tiempo $\tau$ (tau).

\vspace{0.4cm}
\begin{tcolorbox}[colback=gray!5, colframe=gray!50, title=\textbf{Reflexión Inicial EC3}]
Si las leyes de KVL y KCL ahora incluyen ecuaciones que contienen derivadas temporales ($\frac{d}{dt}$), ¿qué tipo de ecuaciones matemáticas gobernarán los circuitos a partir de hoy?
\end{tcolorbox}

\end{document}
